\markdownRendererDocumentBegin
\markdownRendererSectionBegin
\markdownRendererHeadingOne{化学基础知识}\markdownRendererInterblockSeparator
{}\markdownRendererSectionBegin
\markdownRendererHeadingTwo{气体}\markdownRendererInterblockSeparator
{}气体基本特性：\markdownRendererInterblockSeparator
{}\markdownRendererOlBeginTight
\markdownRendererOlItemWithNumber{1}自动扩散性\markdownRendererOlItemEnd 
\markdownRendererOlItemWithNumber{2}可压缩性\markdownRendererOlItemEnd 
\markdownRendererOlEndTight \markdownRendererInterblockSeparator
{}\markdownRendererSectionBegin
\markdownRendererHeadingThree{气体的状态方程}\markdownRendererInterblockSeparator
{}\markdownRendererSectionBegin
\markdownRendererHeadingFour{理想气体的状态方程}\markdownRendererInterblockSeparator
{}假设分子微粒之间的作用力、分子本身的体积忽略不计的气体为\markdownRendererStrongEmphasis{理想气体}。\markdownRendererParagraphSeparator
{}（\markdownRendererStrongEmphasis{低压高温}状态）理想气体状态方程 \markdownRendererDollarSign{}pV=nRT\markdownRendererDollarSign{}\markdownRendererInterblockSeparator
{}
\markdownRendererSectionEnd \markdownRendererSectionBegin
\markdownRendererHeadingFour{实际气体的状态方程}\markdownRendererInterblockSeparator
{}Van der Waals 气体状态方程 \markdownRendererDollarSign{}\markdownRendererBackslash{}left(p+a\markdownRendererBackslash{}left(\markdownRendererLeftBrace{}\markdownRendererBackslash{}dfrac\markdownRendererLeftBrace{}n\markdownRendererRightBrace{}\markdownRendererLeftBrace{}V\markdownRendererRightBrace{}\markdownRendererRightBrace{}\markdownRendererBackslash{}right)\markdownRendererCircumflex{}2\markdownRendererBackslash{}right)(V-bn)=nRT\markdownRendererDollarSign{}，其中 \markdownRendererDollarSign{}p\markdownRendererUnderscore{}内=\markdownRendererBackslash{}left(\markdownRendererBackslash{}dfrac\markdownRendererLeftBrace{}n\markdownRendererUnderscore{}外\markdownRendererRightBrace{}\markdownRendererLeftBrace{}V\markdownRendererRightBrace{}\markdownRendererBackslash{}right)\markdownRendererBackslash{}left(\markdownRendererBackslash{}dfrac\markdownRendererLeftBrace{}n\markdownRendererUnderscore{}内\markdownRendererRightBrace{}\markdownRendererLeftBrace{}V\markdownRendererRightBrace{}\markdownRendererBackslash{}right)=\markdownRendererBackslash{}left(\markdownRendererBackslash{}dfrac\markdownRendererLeftBrace{}n\markdownRendererRightBrace{}\markdownRendererLeftBrace{}V\markdownRendererRightBrace{}\markdownRendererBackslash{}right)\markdownRendererCircumflex{}2\markdownRendererDollarSign{}，\markdownRendererDollarSign{}a\markdownRendererDollarSign{} 是与分子间引力有关的常数，\markdownRendererDollarSign{}b\markdownRendererDollarSign{} 是与分子本身体积有关的常数，统称为 van der Waals 常数。通常，\markdownRendererDollarSign{}a\markdownRendererDollarSign{} 和 \markdownRendererDollarSign{}b\markdownRendererDollarSign{} 数值越大，表明实际气体距理想气体愈远，从而产生的偏差就愈大。\markdownRendererInterblockSeparator
{}
\markdownRendererSectionEnd 
\markdownRendererSectionEnd \markdownRendererSectionBegin
\markdownRendererHeadingThree{混合气体的分压定律}\markdownRendererInterblockSeparator
{}Dalton 分压定律 \markdownRendererDollarSign{}\markdownRendererBackslash{}displaystyle p\markdownRendererUnderscore{}总=\markdownRendererBackslash{}sum\markdownRendererUnderscore{}ip\markdownRendererUnderscore{}i\markdownRendererDollarSign{}\markdownRendererSoftLineBreak
{}\markdownRendererDollarSign{}p\markdownRendererUnderscore{}i=p\markdownRendererUnderscore{}总x\markdownRendererUnderscore{}i\markdownRendererDollarSign{}\markdownRendererParagraphSeparator
{}Amagat 分体积定律 \markdownRendererDollarSign{}\markdownRendererBackslash{}displaystyle V\markdownRendererUnderscore{}总=\markdownRendererBackslash{}sum\markdownRendererUnderscore{}iV\markdownRendererUnderscore{}i\markdownRendererDollarSign{}\markdownRendererInterblockSeparator
{}
\markdownRendererSectionEnd \markdownRendererSectionBegin
\markdownRendererHeadingThree{气体扩散定律}\markdownRendererInterblockSeparator
{}\markdownRendererDollarSign{}v \markdownRendererBackslash{}propto \markdownRendererBackslash{}sqrt\markdownRendererLeftBrace{}\markdownRendererBackslash{}dfrac\markdownRendererLeftBrace{}1\markdownRendererRightBrace{}\markdownRendererLeftBrace{}\markdownRendererBackslash{}rho\markdownRendererRightBrace{}\markdownRendererRightBrace{}\markdownRendererDollarSign{}\markdownRendererSoftLineBreak
{}\markdownRendererDollarSign{}v \markdownRendererBackslash{}propto \markdownRendererBackslash{}sqrt\markdownRendererLeftBrace{}\markdownRendererBackslash{}dfrac\markdownRendererLeftBrace{}1\markdownRendererRightBrace{}\markdownRendererLeftBrace{}M\markdownRendererRightBrace{}\markdownRendererRightBrace{}\markdownRendererDollarSign{}\markdownRendererInterblockSeparator
{}
\markdownRendererSectionEnd \markdownRendererSectionBegin
\markdownRendererHeadingThree{气体分子的速率分布和能量分布}\markdownRendererInterblockSeparator
{}\markdownRendererSectionBegin
\markdownRendererHeadingFour{气体分子的速率分布}\markdownRendererInterblockSeparator
{}1859 年，英国物理学家 Maxwell 用概率论的方法推导出了计算 气体分子运动速率分布的公式，讨论了分子运动速率的分布。后来，奥地利人 L. Boltzmann 用统计力学的方法也得到了相同的结果，从而进一步验证了麦克斯韦的速率分布公式。\markdownRendererInterblockSeparator
{}
\markdownRendererSectionEnd \markdownRendererSectionBegin
\markdownRendererHeadingFour{气体分子的能量分布}\markdownRendererInterblockSeparator
{}在无机化学中，常用下面的近似公式来计算和讨论能量的分布。\markdownRendererParagraphSeparator
{}\markdownRendererDollarSign{}\markdownRendererDollarSign{}\markdownRendererSoftLineBreak
{}f\markdownRendererUnderscore{}\markdownRendererLeftBrace{}E\markdownRendererUnderscore{}0\markdownRendererRightBrace{}=\markdownRendererBackslash{}dfrac\markdownRendererLeftBrace{}N\markdownRendererUnderscore{}i\markdownRendererRightBrace{}\markdownRendererLeftBrace{}N\markdownRendererRightBrace{}=\markdownRendererBackslash{}exp\markdownRendererLeftBrace{}\markdownRendererBackslash{}dfrac\markdownRendererLeftBrace{}-E\markdownRendererUnderscore{}0\markdownRendererRightBrace{}\markdownRendererLeftBrace{}RT\markdownRendererRightBrace{}\markdownRendererRightBrace{}\markdownRendererSoftLineBreak
{}\markdownRendererDollarSign{}\markdownRendererDollarSign{}\markdownRendererInterblockSeparator
{}
\markdownRendererSectionEnd 
\markdownRendererSectionEnd 
\markdownRendererSectionEnd \markdownRendererSectionBegin
\markdownRendererHeadingTwo{液体和溶液}\markdownRendererInterblockSeparator
{}\markdownRendererStrongEmphasis{溶液}是指一种或一种以上的物质以分子或离子的形式溶解在另一种物质中形成的均一 、稳定的混合物。其中物质的量多且能溶解其他物质的组分称为\markdownRendererStrongEmphasis{溶剂}，含量少的组分称为\markdownRendererStrongEmphasis{溶质}。\markdownRendererInterblockSeparator
{}\markdownRendererSectionBegin
\markdownRendererHeadingThree{溶液浓度的表示方法}\markdownRendererInterblockSeparator
{}对于溶剂 \markdownRendererDollarSign{}\markdownRendererBackslash{}rm A\markdownRendererDollarSign{} 和溶质 \markdownRendererDollarSign{}\markdownRendererBackslash{}rm B\markdownRendererDollarSign{}，有\markdownRendererInterblockSeparator
{}\markdownRendererOlBeginTight
\markdownRendererOlItemWithNumber{1}质量摩尔浓度 \markdownRendererDollarSign{}b(\markdownRendererLeftBrace{}\markdownRendererBackslash{}rm B\markdownRendererRightBrace{})=\markdownRendererBackslash{}dfrac\markdownRendererLeftBrace{}n(\markdownRendererLeftBrace{}\markdownRendererBackslash{}rm B\markdownRendererRightBrace{})\markdownRendererRightBrace{}\markdownRendererLeftBrace{}m(\markdownRendererLeftBrace{}\markdownRendererBackslash{}rm A\markdownRendererRightBrace{})\markdownRendererRightBrace{}\markdownRendererDollarSign{}\markdownRendererOlItemEnd 
\markdownRendererOlItemWithNumber{2}物质的量浓度 \markdownRendererDollarSign{}c(\markdownRendererLeftBrace{}\markdownRendererBackslash{}rm B\markdownRendererRightBrace{})=\markdownRendererBackslash{}dfrac\markdownRendererLeftBrace{}n(\markdownRendererLeftBrace{}\markdownRendererBackslash{}rm B\markdownRendererRightBrace{})\markdownRendererRightBrace{}\markdownRendererLeftBrace{}V\markdownRendererRightBrace{}\markdownRendererDollarSign{}\markdownRendererOlItemEnd 
\markdownRendererOlItemWithNumber{3}质量分数 \markdownRendererDollarSign{}w(\markdownRendererLeftBrace{}\markdownRendererBackslash{}rm B\markdownRendererRightBrace{})=\markdownRendererBackslash{}dfrac\markdownRendererLeftBrace{}m(\markdownRendererLeftBrace{}\markdownRendererBackslash{}rm B\markdownRendererRightBrace{})\markdownRendererRightBrace{}\markdownRendererLeftBrace{}m\markdownRendererRightBrace{}\markdownRendererDollarSign{}\markdownRendererOlItemEnd 
\markdownRendererOlItemWithNumber{4}摩尔分数 \markdownRendererDollarSign{}x(\markdownRendererLeftBrace{}\markdownRendererBackslash{}rm B\markdownRendererRightBrace{})=\markdownRendererBackslash{}dfrac\markdownRendererLeftBrace{}n(\markdownRendererLeftBrace{}\markdownRendererBackslash{}rm B\markdownRendererRightBrace{})\markdownRendererRightBrace{}\markdownRendererLeftBrace{}n\markdownRendererRightBrace{}\markdownRendererDollarSign{}\markdownRendererOlItemEnd 
\markdownRendererOlEndTight \markdownRendererInterblockSeparator
{}\markdownRendererDollarSign{}x\markdownRendererBackslash{}approx \markdownRendererBackslash{}dfrac\markdownRendererLeftBrace{}1000\markdownRendererRightBrace{}\markdownRendererLeftBrace{}18\markdownRendererRightBrace{}b\markdownRendererTilde{}\markdownRendererLeftBrace{}\markdownRendererBackslash{}rm mol\markdownRendererCircumflex{}\markdownRendererLeftBrace{}-1\markdownRendererRightBrace{}\markdownRendererRightBrace{}\markdownRendererDollarSign{}\markdownRendererInterblockSeparator
{}
\markdownRendererSectionEnd \markdownRendererSectionBegin
\markdownRendererHeadingThree{溶液的饱和蒸气压}\markdownRendererInterblockSeparator
{}\markdownRendererSectionBegin
\markdownRendererHeadingFour{纯溶剂的饱和蒸气压}\markdownRendererInterblockSeparator
{}当凝聚速率和蒸发速率相等时，饱和蒸气所具有的压强叫做该温度下液体的\markdownRendererStrongEmphasis{饱和蒸气压}，简称\markdownRendererStrongEmphasis{蒸气压}，用 \markdownRendererDollarSign{}p\markdownRendererCircumflex{}*\markdownRendererDollarSign{} 表示。对同一液体来说，其饱和蒸气压随温度的升高而增大。\markdownRendererInterblockSeparator
{}
\markdownRendererSectionEnd \markdownRendererSectionBegin
\markdownRendererHeadingFour{溶液的饱和蒸气压下降}\markdownRendererInterblockSeparator
{}溶液的饱和蒸气压 \markdownRendererDollarSign{}p\markdownRendererDollarSign{} 小于纯溶剂的蒸气压 \markdownRendererDollarSign{}p\markdownRendererCircumflex{}*\markdownRendererDollarSign{}。\markdownRendererSoftLineBreak
{}溶液的浓度越大，溶质的粒子数目越多，溶液的饱和蒸气压比纯溶剂的饱和蒸气压下降得就越多。\markdownRendererParagraphSeparator
{}Raoult 定律 \markdownRendererDollarSign{}p=p\markdownRendererCircumflex{}*x(\markdownRendererLeftBrace{}\markdownRendererBackslash{}rm A\markdownRendererRightBrace{})\markdownRendererDollarSign{}\markdownRendererParagraphSeparator
{}\markdownRendererDollarSign{}\markdownRendererBackslash{}begin\markdownRendererLeftBrace{}align\markdownRendererRightBrace{}\markdownRendererBackslash{}Delta p \markdownRendererAmpersand{}= p\markdownRendererCircumflex{}\markdownRendererEmphasis{-p\markdownRendererBackslash{}\markdownRendererAmpersand{}=p\markdownRendererCircumflex{}}[1-x(\markdownRendererLeftBrace{}\markdownRendererBackslash{}rm A\markdownRendererRightBrace{})]\markdownRendererBackslash{}\markdownRendererAmpersand{}=p\markdownRendererCircumflex{}*x(\markdownRendererLeftBrace{}\markdownRendererBackslash{}rm B\markdownRendererRightBrace{})\markdownRendererBackslash{}end\markdownRendererLeftBrace{}align\markdownRendererRightBrace{}\markdownRendererDollarSign{}\markdownRendererParagraphSeparator
{}稀溶液的饱和蒸气压下降值与稀溶液的质量摩尔浓度成正比。\markdownRendererInterblockSeparator
{}
\markdownRendererSectionEnd 
\markdownRendererSectionEnd \markdownRendererSectionBegin
\markdownRendererHeadingThree{稀溶液的依数性}\markdownRendererInterblockSeparator
{}\markdownRendererLink{稀溶液的饱和蒸气压下降}{\markdownRendererHash{}溶液的饱和蒸气压下降}{#溶液的饱和蒸气压下降}{}是稀溶液的一种依数性质。\markdownRendererInterblockSeparator
{}\markdownRendererSectionBegin
\markdownRendererHeadingFour{溶液的沸点升高}\markdownRendererInterblockSeparator
{}液体汽化的方式有两种，即\markdownRendererStrongEmphasis{蒸发}和\markdownRendererStrongEmphasis{沸腾}。 当液体的饱和蒸气压与外界大气压强相等时，液体的汽化将在其表面和内部同时发生，称为\markdownRendererStrongEmphasis{沸腾}，这时的温度称为该液体的\markdownRendererStrongEmphasis{沸点}。 而在低于此温度下的汽化，仅在液体表面上进行，称为\markdownRendererStrongEmphasis{蒸发}。\markdownRendererParagraphSeparator
{}难挥发性非电解质稀溶液沸点升高值 \markdownRendererDollarSign{}\markdownRendererBackslash{}Delta T\markdownRendererUnderscore{}\markdownRendererBackslash{}mathrm\markdownRendererLeftBrace{}b\markdownRendererRightBrace{}=T\markdownRendererUnderscore{}\markdownRendererBackslash{}mathrm\markdownRendererLeftBrace{}b\markdownRendererRightBrace{}-T\markdownRendererUnderscore{}\markdownRendererBackslash{}mathrm\markdownRendererLeftBrace{}b\markdownRendererRightBrace{}\markdownRendererCircumflex{}*\markdownRendererDollarSign{}，与其饱和蒸气压下降的数值成正比，即 \markdownRendererDollarSign{}\markdownRendererBackslash{}Delta T\markdownRendererUnderscore{}\markdownRendererBackslash{}mathrm\markdownRendererLeftBrace{}b\markdownRendererRightBrace{} \markdownRendererBackslash{}propto \markdownRendererBackslash{}Delta p\markdownRendererDollarSign{}，记 \markdownRendererDollarSign{}k\markdownRendererUnderscore{}\markdownRendererBackslash{}mathrm\markdownRendererLeftBrace{}b\markdownRendererRightBrace{}\markdownRendererDollarSign{} 为稀溶液\markdownRendererStrongEmphasis{沸点升高常数}，得 \markdownRendererDollarSign{}\markdownRendererBackslash{}Delta T\markdownRendererUnderscore{}\markdownRendererBackslash{}mathrm\markdownRendererLeftBrace{}b\markdownRendererRightBrace{}=k\markdownRendererUnderscore{}\markdownRendererBackslash{}mathrm\markdownRendererLeftBrace{}b\markdownRendererRightBrace{}b\markdownRendererDollarSign{}（\markdownRendererStrongEmphasis{沸点升高公式}）\markdownRendererInterblockSeparator
{}
\markdownRendererSectionEnd \markdownRendererSectionBegin
\markdownRendererHeadingFour{溶液的凝固点降低}\markdownRendererInterblockSeparator
{}液体凝固成固体（严格说是晶体）时的温度称为该液体的\markdownRendererStrongEmphasis{凝固点}。由于溶液的蒸气压下降导致溶液的凝固点低于纯溶剂的凝固点。\markdownRendererParagraphSeparator
{}难挥发性非电解质稀溶液凝固点降低值 \markdownRendererDollarSign{}\markdownRendererBackslash{}Delta T\markdownRendererUnderscore{}\markdownRendererBackslash{}mathrm\markdownRendererLeftBrace{}f\markdownRendererRightBrace{}=T\markdownRendererUnderscore{}\markdownRendererBackslash{}mathrm\markdownRendererLeftBrace{}f\markdownRendererRightBrace{}\markdownRendererCircumflex{}*-T\markdownRendererUnderscore{}\markdownRendererBackslash{}mathrm\markdownRendererLeftBrace{}f\markdownRendererRightBrace{}\markdownRendererDollarSign{}，与其饱和蒸气压下降的数值成正比，记 \markdownRendererDollarSign{}k\markdownRendererUnderscore{}\markdownRendererBackslash{}mathrm\markdownRendererLeftBrace{}f\markdownRendererRightBrace{}\markdownRendererDollarSign{} 为\markdownRendererStrongEmphasis{凝固点降低常数}，得 \markdownRendererDollarSign{}\markdownRendererBackslash{}Delta T\markdownRendererUnderscore{}\markdownRendererBackslash{}mathrm\markdownRendererLeftBrace{}f\markdownRendererRightBrace{}=k\markdownRendererUnderscore{}\markdownRendererBackslash{}mathrm\markdownRendererLeftBrace{}f\markdownRendererRightBrace{}b\markdownRendererDollarSign{}（凝固点降低公式）\markdownRendererInterblockSeparator
{}
\markdownRendererSectionEnd \markdownRendererSectionBegin
\markdownRendererHeadingFour{渗透压}\markdownRendererInterblockSeparator
{}溶剂分子透过半透膜进入溶液的现象，称为\markdownRendererStrongEmphasis{渗透现象}。渗透平衡时两侧液面高度差造成的静压称为溶液的\markdownRendererStrongEmphasis{渗透压}。\markdownRendererParagraphSeparator
{}Van’t Hoff 指出 \markdownRendererDollarSign{}\markdownRendererBackslash{}mathit\markdownRendererLeftBrace{}\markdownRendererBackslash{}Pi\markdownRendererRightBrace{}V = nRT\markdownRendererDollarSign{}\markdownRendererInterblockSeparator
{}
\markdownRendererSectionEnd 
\markdownRendererSectionEnd 
\markdownRendererSectionEnd \markdownRendererSectionBegin
\markdownRendererHeadingTwo{固体和晶体}\markdownRendererInterblockSeparator
{}\markdownRendererSectionBegin
\markdownRendererHeadingThree{晶体和非晶体}\markdownRendererInterblockSeparator
{}\markdownRendererSectionBegin
\markdownRendererHeadingFour{非晶体的特性}\markdownRendererInterblockSeparator
{}
\markdownRendererSectionEnd \markdownRendererSectionBegin
\markdownRendererHeadingFour{晶体的特性}\markdownRendererInterblockSeparator
{}
\markdownRendererSectionEnd 
\markdownRendererSectionEnd \markdownRendererSectionBegin
\markdownRendererHeadingThree{对称性}\markdownRendererInterblockSeparator
{}\markdownRendererSectionBegin
\markdownRendererHeadingFour{旋转和旋转轴}\markdownRendererInterblockSeparator
{}
\markdownRendererSectionEnd \markdownRendererSectionBegin
\markdownRendererHeadingFour{反映和镜面}\markdownRendererInterblockSeparator
{}
\markdownRendererSectionEnd \markdownRendererSectionBegin
\markdownRendererHeadingFour{反演和对称中心}\markdownRendererInterblockSeparator
{}
\markdownRendererSectionEnd 
\markdownRendererSectionEnd \markdownRendererSectionBegin
\markdownRendererHeadingThree{晶体和点阵}\markdownRendererInterblockSeparator
{}
\markdownRendererSectionEnd \markdownRendererSectionBegin
\markdownRendererHeadingThree{晶系和点阵型式}\markdownRendererInterblockSeparator
{}\markdownRendererSectionBegin
\markdownRendererHeadingFour{七个晶系}\markdownRendererInterblockSeparator
{}
\markdownRendererSectionEnd \markdownRendererSectionBegin
\markdownRendererHeadingFour{十四种空间点阵型式}\markdownRendererInterblockSeparator
{}
\markdownRendererSectionEnd 
\markdownRendererSectionEnd \markdownRendererSectionBegin
\markdownRendererHeadingThree{晶胞}
\markdownRendererSectionEnd 
\markdownRendererSectionEnd 
\markdownRendererSectionEnd \markdownRendererDocumentEnd